\section{Error Analysis and Noise Characterization}

\subsection{Ideal versus hardware distributions}

Figure~\ref{fig:ideal_hardware_hist} is the wide histogram generated in the executed notebook. It compares a 1024-shot ideal reference with the measured \texttt{ibm\_marrakesh} counts on the same selected support. The standard distribution is moderately distorted. The \GMQAOA{} distribution exhibits a much larger collapse of the dominant ideal peaks and substantial population of states having zero ideal probability.

\begin{figure}[H]
\centering
\includegraphics[width=0.99\textwidth]{ideal_vs_ibm_marrakesh_histograms.png}
\caption{Ideal finite-shot reference versus IBM Quantum hardware for standard \QAOA{} and \GMQAOA{}. The bottom panel makes the collapse of the ideal \GMQAOA{} peaks visible.}
\label{fig:ideal_hardware_hist}
\end{figure}

\subsection{Feasible and optimal probabilities}

Table~\ref{tab:probability_degradation} compares exact ideal values with complete hardware counts. The standard ideal values and its global distribution metrics are reconstructed within the six-decimal precision of the repository exports, as explained in Section~6.

\begin{table}[H]
\centering
\caption{Ideal-to-hardware probability degradation.}
\label{tab:probability_degradation}
\small
\begin{tabular}{lrrrrrr}
\toprule
Algorithm & Ideal $\PF$ & QPU $\PF$ & Ideal $\Popt$ & QPU $\Popt$ & $\Delta\Popt$ & Relative $\Popt$ error\\
\midrule
Standard \QAOA{} & 0.091582 & 0.077148 & 0.023380 & 0.015625 & $-0.007755$ & $-33.17\%$\\
\GMQAOA{} & 1.000000 & 0.119141 & 0.503666 & 0.028320 & $-0.475346$ & $-94.38\%$\\
\bottomrule
\end{tabular}
\end{table}

\begin{figure}[H]
\centering
\includegraphics[width=0.96\textwidth]{ideal_hardware_quality_comparison.png}
\caption{Ideal and physical feasible/optimal probabilities for the hardware instance.}
\label{fig:quality_comparison}
\end{figure}

The 95\% Wilson intervals are
\begin{align}
P_{\mathrm{F,QPU}}^{\mathrm{std}}&\in[0.06234,0.09512],&
P_{\mathrm{opt,QPU}}^{\mathrm{std}}&\in[0.00964,0.02523],\\
P_{\mathrm{F,QPU}}^{\mathrm{GM}}&\in[0.10071,0.14042],&
P_{\mathrm{opt,QPU}}^{\mathrm{GM}}&\in[0.01979,0.04038].
\end{align}
The standard ideal optimum probability lies inside its physical interval, whereas the \GMQAOA{} ideal optimum probability is far outside the physical interval.

\subsection{Global distribution metrics}

\begin{table}[H]
\centering
\caption{Distribution-level errors. Standard values marked $\dagger$ are reconstructed within the precision of the exported ideal probabilities.}
\label{tab:distribution_metrics}
\begin{tabular}{lrrrr}
\toprule
Algorithm & TVD & Hellinger & $\Fcl$ & $1-\TVD$\\
\midrule
Standard \QAOA{}$^\dagger$ & 0.2082 & 0.2346 & 0.8929 & 0.7918\\
\GMQAOA{} & 0.8809 & 0.8284 & 0.0984 & 0.1191\\
\bottomrule
\end{tabular}
\end{table}

\begin{figure}[H]
\centering
\includegraphics[width=0.76\textwidth]{distribution_error_metrics.png}
\caption{Distribution-level distance and infidelity measures derived from the complete archived counts.}
\label{fig:error_metrics}
\end{figure}

For \GMQAOA{}, the ideal distribution has support only on $\mathcal F$. All six measured feasible probabilities remain below their ideal values, so the TVD equals the measured leakage:
\begin{equation}
\TVD_{\mathrm{GM}}=1-P_{\mathrm{F,QPU}}^{\mathrm{GM}}=0.880859375.
\end{equation}
This equality provides a direct interpretation of the error: $88.09\%$ of measured probability has leaked outside the ideal support.

\subsection{Energy distortion}

The expected penalized energies change as follows:
\begin{table}[H]
\centering
\caption{Ideal and observed expected \QUBO{} energies.}
\label{tab:energy_distortion}
\begin{tabular}{lrrr}
\toprule
Algorithm & Ideal $\langle Q\rangle$ & QPU $\langle Q\rangle$ & Difference\\
\midrule
Standard \QAOA{}$^\dagger$ & 18.3443 & 19.6276 & $+1.2833$\\
\GMQAOA{} & $-2.8453$ & 20.2492 & $+23.0945$\\
\bottomrule
\end{tabular}
\end{table}
The \GMQAOA{} expected energy loses the advantage implied by its ideal feasible-state concentration, even though one optimal sample remains visible.

\subsection{Noise mechanisms and evidence}

The experiment does not isolate a single microscopic error channel. The observed distortion is consistent with the combined effects of:
\begin{itemize}
  \item synthesis of controlled state preparations and the multi-controlled phase;
  \item routing on a sparse coupling graph;
  \item accumulated CZ and single-qubit gate error;
  \item decoherence during a substantially longer schedule;
  \item measurement error;
  \item calibration drift and stochastic variation across twirling realizations.
\end{itemize}
Dynamical decoupling and gate twirling were enabled, so the results already include those suppression techniques. Their presence does not constitute a no-mitigation control experiment. The resource contrast strongly supports circuit length as a major explanatory factor, but it does not prove that depth alone caused the observed difference.

\subsection{Best-sample error versus statistical reliability}

Both algorithms have zero best-sample optimization gap because the optimum appeared at least once. This statement is compatible with low reliability: the optimal state occurred only 16 times for standard \QAOA{} and 29 times for \GMQAOA{} out of 1024 shots. A publication claim should therefore report the gap jointly with $\Popt$, confidence intervals, $\PF$, and global distances.
